\chapter{Conclusion and Future Work}
\label{chap:conclusion}
\begingroup
\justifying
\setlength{\parindent}{2em}
\setlength{\parskip}{0.55\baselineskip}

\section{Conclusions}

I set out to run two released RoboBenchMart checkpoints through the benchmark's atomic protocol, using nothing but what is publicly available, and to report what happened. The short answer is that the protocol runs, the results are far from the published ones, and the most likely reason lies in the gap between the checkpoints and the environment they are now evaluated in rather than in the policies as such.

On RQ1, all 42 configurations completed for both Octo and $\pi_{0.5}$, giving 2,520 episodes. Three repairs were needed: an import that an upstream change had removed from the floor environment, the way the client reads instructions through gymnasium wrappers, and per-connection policy state in the Octo server. The evidence that remains is a success flag for each episode, linked to its seed, its log line and its video, with hashes for every file. Per-step trajectories were not saved, because the evaluation script does not ask for them.

On RQ2, the pick-and-place families produced no successes at all, 0 of 900 episodes for each model. The door tasks are the only place where the models differ. $\pi_{0.5}$ closed a door in 37 of 180 episodes and opened one in 2 of 180; Octo closed a door in 4 of 180 and never opened one. Within $\pi_{0.5}$'s closing results, Seeds, Pose and Layout (15, 10 and 12 of 60) are not distinguishable at this sample size. Pairing could not show anything, because the tasks it applies to were already at zero.

On RQ3, the distance from the upstream report is large and systematic. Upstream reports, for example, 82--89\% for $\pi_{0.5}$ on Basket with training seeds, where I observed 0 of 90. Sampling variation cannot produce a gap like that. Since the checkpoints were published, the upstream code has changed the robot's initial joint configuration, the lighting, the robot class and the motion planner, and the demonstration data and scenes have been regenerated; the evaluated release dates from January 2026, the checkpoints from November 2025. I consider this version drift the leading explanation. I have not tested it, and it remains a hypothesis. Several properties of the protocol also bound interpretation regardless of version: door Layout reuses training scenes, every episode starts close to its fixture, actions are executed in open-loop chunks with success read only at chunk boundaries, and the disturbance check covers the whole store.

In plainer terms, the released checkpoints, run on the current benchmark, do not reproduce the published results, and with the records available it is not possible to say exactly why. That is a less satisfying conclusion than a leaderboard. I think it is the more useful one for anyone who intends to use these checkpoints as baselines.

\section{Limitations}

The sample size is the most immediate limitation. Thirty episodes per configuration were enough to establish that pick-and-place rates are low and that $\pi_{0.5}$ closes doors far more often than it opens them, but not enough to resolve differences between conditions of the size the protocol was designed to detect.

The study ran each protocol once. Execution differed between the two models (Octo partly serial and partly on four servers, $\pi_{0.5}$ on one shared server), so the comparison between them is seed-matched but not strictly paired, and run-to-run variation is unmeasured.

The evidence is outcome-level. Without per-step trajectories, and without anyone having watched the videos systematically, I cannot say at which stage failed episodes went wrong. I have deliberately not speculated about failure causes, and the upstream failure taxonomy remains unapplied.

Several facts about the setup are inferred from code rather than observed at run time: the mapping of the locked action indices to the head joints, and the $\pi_{0.5}$ action-chunk length, which the adapter inherits from an OpenPI default I did not check on the revision used. The GPU and several library versions were not recorded during the runs.

Finally, the scope is narrow by design: two checkpoints, atomic tasks only, simulation only. Nothing here speaks to $\pi_0$ or SmolVLA, to composite tasks, or to physical robots.

\section{Future work}
\label{sec:future}

The first experiment I would run is a \textbf{version-aligned re-evaluation}. Check out the upstream code as of early November 2025, pair it with the scene release that preceded the January regeneration (\texttt{demo\_envs} at \texttt{2fd1681a}), and evaluate both checkpoints on a small number of configurations, perhaps two to four, chosen where the upstream and current results differ most, such as $\pi_{0.5}$ Basket and Board under Seeds. Thirty episodes per configuration would be enough. If the rates move close to the upstream figures, the drift hypothesis is confirmed and the individual changes can then be reintroduced one at a time, starting with the initial joint configuration. If they do not, the interface between checkpoint and environment, including action-dimension order and normalisation statistics, is the next place to look.

The second is an \textbf{execution-horizon sweep}. Once a configuration with a non-trivial success rate is available, executing only the first $H \in \{1, 5, 10, 25, 50\}$ actions of each chunk before re-querying the policy would show how much the open-loop execution costs. This needs no retraining. It would also be worth checking the success predicate after every step rather than only at chunk boundaries, which isolates a separate effect.

The third is \textbf{better records}. Running with \texttt{--save-traj} for a subset of configurations would make per-step states available for analysis. With those and the existing videos, the upstream failure taxonomy could be applied under a written coding manual, ideally by more than one annotator, so that failure stages can be reported rather than guessed. Printing the controller's joint list and the policy's action horizon at the start of each run would remove the two remaining inferences.

Only after these steps would it make sense to increase the number of episodes per configuration and revisit the questions the four-condition ladder was designed to answer: how much performance falls from training seeds to new starting poses, to new layouts, and to new task--item pairings.

\section{A staged route to physical validation}

Simulation results, even reproducible ones, are a filter rather than a destination. If a version-aligned baseline is established, a sensible route to physical validation would start with a single shelving unit matching the simulated dimensions, with calibrated cameras placed as in simulation, to separate visual domain shift from contact and controller effects. Real packaging, which is glossy, deformable or transparent in ways the rigid-body simulator does not capture, would come next. Only then would store-scale navigation and safety envelopes around people be worth adding, and reporting would shift from binary success towards measures a store operator cares about, such as items handled per hour, damage to neighbouring stock and the frequency of human intervention.

The step that comes before all of this is less glamorous. A released checkpoint needs a matching, versioned evaluation environment, and the version information has to be recorded at the time the results are produced. The two runs in this dissertation are an attempt to keep such a record.

\endgroup
